\clearpage
\phantomsection
\addcontentsline{toc}{chapter}{原课程信息与课历}
\markboth{原课程信息与课历}{}
\noindent{\color{structurecolor}\bfseries 原课程信息与课历}\par

\noindent\textbf{课程身份与先修.} MIT 18.03 Differential Equations, Spring 2010, 为数学系本科课程, 原课教师为 Haynes Miller 与 Arthur Mattuck. 每周三次一小时讲课、两次一小时习题课. 先修18.01一元微积分, 同修18.02多元微积分; 允许与18.03同时学习18.02. 本册的矩阵记号与必要计算在正文中说明.

\noindent\textbf{原课程教材.} C. Edwards 与 D. Penney, \textit{Elementary Differential Equations with Boundary Value Problems}, 第6版, Prentice Hall, 2003, ISBN 9780136006138; 大纲亦允许第5版, ISBN 9780131457744. 原课还使用 Arthur Mattuck 的 \textit{18.03 Notes and Exercises} 和 Haynes Miller 的 \textit{18.03 Supplementary Notes}. 后文的 EP 表示该商业教材, Notes 表示 Mattuck 讲义中的题号, SN 表示 Miller 补充讲义. 商业教材的题号定位不等于公开题面, 本册没有按猜测复原它们.

\noindent\textbf{原课程大纲.} 原课学习常微分方程及物理系统建模: 一阶方程的解析、图解与数值方法; 高阶线性方程, 尤其常系数二阶方程; 待定系数与常数变易; 复数、复指数以及正弦和指数输入; 振动、阻尼、共振与频率响应; Fourier 级数及周期解; 阶跃、脉冲、卷积与 Laplace 方法; 矩阵、特征值、特征向量与一阶系统; 非线性自治系统的平衡分析及相平面. 本册第1--13章沿数学依赖重组该主线, 并补入已注明的理论内容.

\noindent\textbf{教学组织.} 原课堂通过构造、求解和解释方程训练主动参与; 学生应准备积极参加课堂. 第一次习题课发给投票卡片, 此后每次讲课携带, 用来对课堂问题投票; 意见不一致时讨论, 课末可有一分钟短反馈. 习题课以小组讨论、解题和答疑为主, 鼓励尽早、经常提问, 助教另有办公时间. 原课另设由有经验的本科生值班的辅导室, 可在此完成作业, 考前增派人员. 互动数学小程序 Mathlets 在课堂及每次作业中使用; 下文保留原题所要求的观察任务, 以解析论证和本册独立图形说明数学结论, 不伪称运行过2010年的 Java 小程序.

\noindent\textbf{作业与考核.} 九次作业均分两部分: 第一部分指定教材或 Notes 题目, 第二部分提供更具综合性的公开题面. 应随讲次练习, 而非截止日前一次完成. 原作业说明允许合作讨论, 要求各自独立写出答案并在首页列合作者姓名; 因截止后立即公布答案, 原说明不接受延期. 常规作业按100原始分计, 四讲各24分, 整洁与表达4分; 作业6、7对应三讲, 每讲32分, 表达4分. 学期总评采用下表的885分权重, 与单次作业的100分原始尺度区分.
\begin{center}
\begin{tabular}{lr}\toprule 项目 & 总评点数 \\\midrule
九次作业 & 225 \\
三次一小时期中考试 & 300 \\
一次三小时综合期末考试 & 360 \\\midrule
合计 & 885 \\\bottomrule
\end{tabular}
\end{center}
当期期中试卷题头日期分别为2010年2月24日、3月17日、4月23日. 当期期末题头为2010年5月18日9:00--12:00. 期末没有公开官方答案. 官网练习期末题头写2010, 所链接答案题头写 Spring 2008; 两者若有内容差异, 本册在相应解答中分别说明, 不仅凭链接位置配对.

\noindent\textbf{原课十项核心技能.}
原大纲要求学生个人熟练掌握这些技能, 因其会在其他 MIT 课程中使用; 技能清单向以18.03为先修的任课教师广泛提供. 大纲当时记载140门课程将18.03列为先修或同修, 此数字是2010课程大纲的历史陈述, 不是当前课程数量统计.
\begin{enumerate}
\item 建模得到一阶方程, 用方向场与等斜线理解解, 用 Euler 方法近似.
\item 用积分因子或常数变易求解一阶线性方程.
\item 熟练运算复数和复指数.
\item 求解输入为指数乘多项式的常系数二阶初值问题; 正弦输入时求增益及相位差.
\item 计算 Fourier 系数, 以级数求线性方程的周期解.
\item 用脉冲描述突变, 求单位脉冲响应, 以卷积表示一般输入响应.
\item 使用 Laplace 表求权函数或单位脉冲响应及线性初值解, 联系极点、阻尼与频率响应.
\item 求特征值、特征向量与矩阵指数, 解一阶系统并联系高阶方程.
\item 从迹与行列式重建二维线性自治系统相图.
\item 通过平衡点附近分析判断二维非线性自治系统的定性行为.
\end{enumerate}

\noindent\textbf{公开题面的编排与解答身份.} 正文之后收录七个 Notes 题库单元的217个主编号、九次作业的公开自带题面、22次有公开题面的习题课、三次当期期中与一期末及四套练习考. 第一部分引用 Notes 者按题号及指定小问回指题库, 同一题不再算一个新题; EP 中未公开题面者只保留中文作业指定和准确定位. 每题保留原编号和原条件, 解答标明官方答案译审、译审并补证或编者补解. 编者补解为 AI 辅助形成的独立解答, 与官方答案分开标注, 由另一独立第二读者核算; 同模型实例的交叉审阅不冒充外部专家审稿. 原题或原答的可证实错漏另注明, 不静默改题. 末级小问、出处、重复回指及补解状态见随书逐题清单.

\noindent\textbf{原课程课历.} 以下完整保留官网的讲课 L、习题课 R 顺序、主题、首次引入的技能及作业发布/提交节点. 官网课历没有逐行给出公历日期; 表中不补造日期. 期中I、III的官网行未标 L 编号; L19明确是期中II. R6、R12、R20与R26是复习节点, 无单独公开习题文件; 其余22份题面实际存在, 包括 R23.
\begin{longtable}{p{.065\textwidth}p{.22\textwidth}p{.47\textwidth}p{.105\textwidth}}
\toprule 节点 & 主题 & 首次引入的技能与概念 & 作业与考试 \\\midrule\endhead
\multicolumn{4}{l}{\textbf{一、一阶微分方程}}\\
R1 & 自然增长与分离变量 & 指数增长及捕获建模; 增长率; 分离变量; 通解与特解; 合并积分常数; $\ln|y|$ 及消去; 补回损失解; 用常数满足初值 & -- \\
L1 & 方向场与存在唯一性 & 方向场、积分曲线、等斜线、漏斗; 隐式解; 导数无穷导致不能延续 & 发布作业1 \\
R2 & 方向场与分界曲线 & 分离轨线及解的极值 & -- \\
L2 & 数值方法 & Euler 方法 & -- \\
L3 & 线性方程与建模 & 一阶线性方程; 系统与信号视角; 银行账户; RC 电路; 常输入的分离变量解 & -- \\
R3 & Euler 方法及线性模型 & 混盐问题 & -- \\
L4 & 积分因子 & 齐次方程与零输入; 积分因子; 暂态; 扩散模型及耦合常数 & -- \\
R4 & 一阶线性方程 & 正弦输入 & -- \\
L5 & 复数与单位根 & 复数; 单位根 & 提交1; 发布2 \\
L6 & 复指数与正弦函数 & 复指数; 正弦的振幅、角频率及相位滞后 & -- \\
L7 & 指数及正弦输入的线性响应 & 一阶线性系统的指数/正弦响应; 正弦输入的复方程; 作业中的半衰期 & -- \\
R5 & 复数与复指数 & 原栏为空 & -- \\
L8 & 自治方程、相线与稳定性 & 自治方程; 相线; 稳定性; $e^{k(t-t_0)}$ 与 $ce^{kt}$ & 提交2; 发布3 \\
L9 & 线性与非线性 & 解不能延续的情形 & -- \\
R6 & 期中I复习 & 原栏为空 & 无独立题面 \\
-- & 期中I & 原栏为空 & 一小时考试 \\
\multicolumn{4}{l}{\textbf{二、二阶线性方程}}\\
R7 & 二阶方程的解 & 简谐振子; 初值; 齐次叠加 & -- \\
L11 & 模态与特征多项式 & 弹簧、质量及阻尼器; 一般二阶线性方程; 特征多项式; 实根解 & -- \\
L12 & 振动与阻尼 & 复根; 欠/过/临界阻尼; 复数替换与实解提取; 暂态; 根图 & -- \\
R8 & 二阶常系数齐次方程 & 一般正弦响应; 归一化解 & -- \\
L13 & 指数响应与弹簧驱动 & 受迫系统; 叠加; 指数响应公式; 复数替换; 正弦输入的正弦响应 & -- \\
R9 & 指数与正弦输入 & 原栏为空 & -- \\
L14 & 复增益与阻尼器驱动 & 增益、相位滞后及复增益 & 提交3; 发布4 \\
L15 & 算子、待定系数与共振 & 微分算子; 共振; 待定系数 & -- \\
R10 & 增益、共振与待定系数 & 原栏为空 & -- \\
L16 & 频率响应 & 频率响应 & -- \\
R11 & 频率响应 & 一阶频率响应 & -- \\
L17 & LTI 系统与 RLC 电路 & RLC 电路; 时间平移不变性 & 提交4; 发布5 \\
L18 & 工程应用 & 阻尼比 & -- \\
R12 & 期中II复习 & 原栏为空 & 无独立题面 \\
L19 & 期中II & 原栏为空 & 一小时考试 \\
\multicolumn{4}{l}{\textbf{三、Fourier 级数及变换响应}}\\
R13 & Fourier 级数入门 & 周期函数 & -- \\
L20 & Fourier 级数 & Fourier 级数; 正交性; Fourier 积分 & -- \\
L21 & Fourier 级数的运算 & 方波; 分段连续; 三角恒等式、线性组合及平移 & -- \\
R14 & Fourier 级数 & 不同周期 & -- \\
L22 & 周期解与共振 & 级数求导与积分; 谐波响应; Fourier 级数的振幅--相位表达 & -- \\
R15 & Fourier 谐波响应 & 原栏为空 & -- \\
L23 & 阶跃及脉冲 & 阶跃; $\delta$; 正则函数与奇异函数; 广义函数及导数 & 提交5; 发布6 \\
L24 & 阶跃响应与脉冲响应 & 单位阶跃/脉冲响应; 静止初值; 一阶、二阶响应 & -- \\
R16 & 阶跃、脉冲及其响应 & 原栏为空 & -- \\
L25 & 卷积 & 脉冲后的初值; 时间不变与 $D$、时间平移的交换; 卷积乘积; 初值响应 $w*q$ & -- \\
R17 & 卷积 & $\delta$ 为卷积单位元 & -- \\
L26 & Laplace 基本性质 & Laplace 变换; 收敛域; $\Lap[t^n]$; $s$ 平移; 正弦、余弦变换; 时域与变换域 & 提交6; 发布7 \\
L27 & Laplace 解方程 & $\Lap[\delta]$; 时间导数规则; 逆变换; 部分分式与遮盖法; 一阶非静止初值 & -- \\
R18 & Laplace 变换 & 用变换求单位阶跃响应 & -- \\
L28 & 二阶方程与配方 & $s$ 求导规则; 二阶方程 & -- \\
R19 & Laplace 方法II & 原栏为空 & -- \\
L29 & 极点图 & 权函数、传递函数及其变换关系; 时间平移; 极点及长期行为 & 提交7; 发布8 \\
L30 & 传递函数与频率响应 & 稳定性; 传递函数与增益 & -- \\
R20 & 期中III复习 & 原栏为空 & 无独立题面 \\
-- & 期中III & 原栏为空 & 一小时考试 \\
\multicolumn{4}{l}{\textbf{四、一阶方程组}}\\
L32 & 线性系统与矩阵 & 一阶线性系统; 消元; 矩阵; 反消元及伴随矩阵 & -- \\
R21 & 一阶线性系统 & 原栏为空 & -- \\
L33 & 特征值与特征向量 & 行列式; 特征值; 特征向量; 初值 & -- \\
R22 & 特征值与特征向量 & 解与轨道的区别 & -- \\
L34 & 复特征值或重根 & 特征值与系数; 复特征值; 重根; 缺陷情形与完整特征基 & 提交8; 发布9 \\
L35 & 线性系统相平面 & 迹--行列式平面; 稳定性 & -- \\
R23 & 线性相图 & 线性相图随参数变形 & -- \\
L36 & 正常模态与矩阵指数 & 矩阵指数; 解耦系统; 指数律 & -- \\
R24 & 矩阵指数 & 非齐次系统, 常输入 & -- \\
L37 & 非线性系统 & 非线性自治系统; 向量场; 相图; 平衡点; 平衡附近线性化; Jacobian 矩阵 & 提交9 \\
L38 & 平衡附近线性化及非线性单摆 & 非线性单摆; 飞机长周期振荡; Tacoma Narrows 桥 & -- \\
R25 & 自治系统 & 捕食--被捕食系统 & -- \\
L39 & 线性方法的局限 & 结构稳定; 极限环; 奇异吸引子 & -- \\
R26 & 课程复习 & 原栏为空 & 无独立题面 \\
-- & 综合期末 & 原栏为空 & 三小时考试 \\
\bottomrule
\end{longtable}


来源: \href{https://ocw.mit.edu/courses/18-03-differential-equations-spring-2010/pages/syllabus/}{原课程大纲}, \href{https://ocw.mit.edu/courses/18-03-differential-equations-spring-2010/pages/calendar/}{原课程课历}, \href{https://ocw.mit.edu/courses/18-03-differential-equations-spring-2010/pages/assignments/}{作业目录}, \href{https://ocw.mit.edu/courses/18-03-differential-equations-spring-2010/pages/recitations/}{习题课目录}, \href{https://ocw.mit.edu/courses/18-03-differential-equations-spring-2010/pages/exams/}{考试目录}. 课程信息仅说明2010年的实际安排, 不是2027年招生考纲的认定.
